\documentclass{article}


% if you need to pass options to natbib, use, e.g.:
\PassOptionsToPackage{numbers, compress}{natbib}
% before loading neurips_2023


% ready for submission
% \usepackage{neurips_2023}


% to compile a preprint version, e.g., for submission to arXiv, add add the
% [preprint] option:
\usepackage[preprint]{neurips_2023}


% to compile a camera-ready version, add the [final] option, e.g.:
%     \usepackage[final]{neurips_2023}


% to avoid loading the natbib package, add option nonatbib:
   % \usepackage[nonatbib]{neurips_2023}


\usepackage[utf8]{inputenc} % allow utf-8 input
\usepackage[T1]{fontenc}    % use 8-bit T1 fonts
\usepackage{hyperref}       % hyperlinks
\hypersetup{                %% hyperlinks settings
    colorlinks=true,
    linkcolor=blue,
    filecolor=blue,      
    urlcolor=blue,
    citecolor=cyan, %% green
    pdftitle={My document},
    pdfpagemode=FullScreen,
    pdfborder={0 0 0},
}
\usepackage{url}            % simple URL typesetting
\usepackage{booktabs}       % professional-quality tables
\usepackage{amsfonts}       % blackboard math symbols
\usepackage{nicefrac}       % compact symbols for 1/2, etc.
\usepackage{microtype}      % microtypography
\usepackage{xcolor}         % colors
\usepackage{graphicx}       %% pdf picture
\usepackage{wrapfig}        %% text around figure
\usepackage{amsmath}        %% math
\usepackage{amssymb}
\usepackage{multirow}       %% table
\usepackage{multicol}
\usepackage{algorithm}      %% algotithm
\usepackage{algpseudocode}

% \newcommand{\zj}[1]{\textcolor{red}{#1}}

\title{DiffSketcher: Text Guided Vector Sketch Synthesis through Latent Diffusion Models}


% The \author macro works with any number of authors. There are two commands
% used to separate the names and addresses of multiple authors: \And and \AND.
%
% Using \And between authors leaves it to LaTeX to determine where to break the
% lines. Using \AND forces a line break at that point. So, if LaTeX puts 3 of 4
% authors names on the first line, and the last on the second line, try using
% \AND instead of \And before the third author name.

% \author{
% Preprint. Under Review.
% }

\author{%
  Ximing Xing \\
  Beihang University\\
  \texttt{ximingxing@buaa.edu.cn} \\
  \And
  Chuang Wang \\
  Beihang University\\
  \texttt{18377011@buaa.edu.cn} \\
  \And
  Haitao Zhou \\
  Beihang University\\
  \texttt{18377221@buaa.edu.cn} \\
  \And
  Jing Zhang \\
  Beihang University\\
  \texttt{zhang\_jing@buaa.edu.cn} \\
  \And
  Qian Yu\thanks{Corresponding author} \\
  Beihang University\\
  \texttt{qianyu@buaa.edu.cn} \\
  \And
  Dong Xu \\
  HongKong University\\
  \texttt{dongxu@cs.hku.hk} \\
}


\begin{document}


\maketitle

\begin{figure*}[h]
\centering
\includegraphics[width=1\linewidth]{teaser.pdf}
\vspace{-1.5em}
\caption{
\textbf{Top}: Visualizations of the vector sketches generated by our proposed method, \textit{DiffSketcher}. \textbf{Bottom}: Visualizations of the drawing process. For each example, we show two sketches with a different number of strokes. 
}
\label{fig:tears}
\end{figure*}
\begin{abstract}
% Sketching is a powerful method of expressing visual information due to its abstract and minimalistic nature. However, sketching requires a semantic understanding and prior knowledge of high-level concepts. 
%Minimal representations, such as sketches and natural language, are powerful tools for effectively communicating ideas by emphasizing their essence. 
%While there have been significant advancements in text-to-image generation, the task of text-to-sketch generation remains largely unexplored.
%The combination of these two fundamental forms of human expression is crucial for effective communication, although translating from text to sketches is more challenging than captioning.
Even though trained mainly on images, we discover that pretrained diffusion models show impressive power in guiding sketch synthesis. 
In this paper, we present DiffSketcher, an innovative algorithm that creates \textit{vectorized} free-hand sketches using natural language input. 
DiffSketcher is developed based on a pre-trained text-to-image diffusion model. 
%It improves the sketch's fidelity and consistency with the input text by 
It performs the task by directly optimizing a set of Bézier curves with an extended version of the score distillation sampling (SDS) loss, which allows us to use a raster-level diffusion model as a prior for optimizing a parametric vectorized sketch generator. Furthermore, we explore attention maps embedded in the diffusion model for effective stroke initialization to speed up the generation process.
%We optimize an attentive-initialized vectorized sketch model via gradient descent such that its raster renderings from attention prior achieve a low loss. 
The generated sketches demonstrate multiple levels of abstraction while maintaining recognizability, underlying structure, and essential visual details of the subject drawn. 
Our experiments show that DiffSketcher achieves greater quality than prior work. The code and demo of DiffSketcher can be found at \href{https://ximinng.github.io/DiffSketcher-project/}{this link}.
\end{abstract}

% Vector graphics can be scaled to any size, and are compact.
\section{Introduction}
Minimal representations, such as sketches and natural language, are powerful tools for effectively conveying ideas by emphasizing the subject's essence. As the carriers of abstract concepts, natural language conveys abstract semantic understanding, whereas sketches embody the human visual abstract presentation. A sketch can provide more visual details than natural language, making the concept more concrete. % provides a practical example. 
When a designer discusses the design plan with a client, the designer may sketch a prototype based on the client's description to ensure a full understanding of his/her requirements. If the process of text-to-sketch could be learned automatically, it would significantly lessen associated labor costs.


%while sketches correspond to the human visual abstract presentation while natural language to abstract semantic understanding. 
% In some cases, natural language can briefly express the subject's essence but lacks concreteness. Given a description, people may have different imaginations. As the visual abstraction, a sketch provides more visual details which makes the concept more concrete. Therefore, it is 
%The organic combination of these two fundamental forms of human expression is crucial for effective communication of ideas. While the ability to translate from sketches to textual descriptions is called captioning, translating from text to sketches is more challenging.

Unfortunately, the task of text-to-sketch remains unexplored. Some studies have been conducted on generating image-conditioned sketches~\cite{infodrawing_Chan_2022,liu2020unsupervised,vinker_2022_clipasso,clipascene_vinker_2022}. One such method is Info-drawing~\cite{infodrawing_Chan_2022}, which treats a sketch as a photo and uses a generative adversarial network (GAN) to generate a raster sketch based on an input image. It also introduces a style loss to ensure the generated sketches have a similar style to the reference sketch used in training. CLIPasso~\cite{vinker_2022_clipasso} presents a novel pipeline to produce vector sketches through a differentiable rasterizer optimized by a semantic loss. CLIPasso can only handle the images with a single object, so its follow-up work CLIPascene~\cite{clipascene_vinker_2022} extends to scene-level images. However, all of these methods depend on input images and do not facilitate the generation of sketches directly from text inputs. 
Moreover, despite their ability to generate realistic sketches from photos, these methods have a limitation in that they cannot generate \textit{new content}.

Recent breakthroughs in text-to-image generation have been driven by diffusion models~\cite{GLIDE_2022_nichol,DALLE2_2022_ramesh,ldm_2022_Rombach,imagen_2022_saharia} trained on billions of image-text pairs~\cite{laion_schuhmann_2022}. These models, conditioned on text, now support high-fidelity, diverse, and controllable image synthesis. Nonetheless, the present text-to-image diffusion models cannot produce highly abstract and vectorized free-hand sketches (as shown in Fig.~\ref{fig:more_results}).
%Unfortunately, the research on text-to-sketch generation remains unexplored.
We draw inspiration from the highly effective text-to-image diffusion models and image-conditioned sketch generators to build a bridge between two fundamental forms of human expression, namely text and free-hand sketches, resulting in the development of our proposed text-to-sketch generator.
% This motivates us to bridge the two primitive abstract expressions of human beings, i.e., text and free-hand sketches, by leveraging the power of the pretrained text-to-image diffusion models.
%in order to achieve effective expression of ideas.

In this work, we present DiffSketcher, an algorithm that can synthesize novel and high-quality free-hand vector sketches based on natural language text input. DiffSketcher does not require any text-to-sketch training pairs or large datasets of sketches. Instead, with the guidance of a pretrained diffusion model~\cite{ldm_2022_Rombach}, we define a sketch as a set of Bézier curves and optimize the curves' parameters through a differentiable rasterizer~\cite{Li_2020_diffvg}. The key idea behind DiffSketcher is to transfer the prior knowledge from text-to-image generation model into the differentiable rasterizer. This allows the optimization of text-guided sketches with semantic consistency, ensuring that the final output sketches are coherent and aligned with their corresponding text inputs.
However, it is a non-trivial task to fully leverage the pre-trained diffusion model to efficiently generate high-quality free-hand sketches for both simple objects and complex scenes.
%through an extended SDS loss inspired from DreamFusion~\cite{poole_2023_dreamfusion}
% In this work, we present DiffSketcher, an algorithm that can synthesize novel and high-quality free-hand vector sketches based on natural language input. DiffSketcher does not require any training or large datasets of sketches. Instead, we define a sketch as a set of Bézier curves and optimize the curves' parameters through a differentiable rasterizer~\cite{Li_2020_diffvg}, which is guided by a pretrained diffusion model~\cite{ldm_2022_Rombach}. The key idea is to transfer image synthesis prior knowledge into the differentiable rasterizer, with the similar idea of DreamFusion~\cite{poole_2023_dreamfusion}. 
% \zj{Surprisingly, even a perceptual loss can coordinate these two parts and generate reasonable sketches. Specifically, our algorithm samples an image from the diffusion model, vectorizes it automatically, and outputs a high-quality sketch.}
% Surprisingly, even a perceptual loss can coordinate these two parts and generate reasonable sketches. To be specific, sampling an image from the diffusion model, then vectorizing it automatically.

To this end, we propose three strategies to improve the generation quality and efficiency: 
(1) Present an extended Score Distillation Sampling (SDS) Loss to guide the optimization of curve parameters.
% (2) Compare the effect of different loss functions, including CLIP loss, LPIPS loss, and Score Distill Sampling (SDS) Loss. 
%%%We first optimize the differentiable rasterizer with the CLIP loss and LPIPS loss. The gradient is computed by comparing the similarity between the photo generated by the diffusion model and the sketch rendered by the rasterizer. 
%LPIPS loss is widely used in style transfer, comparing the perceptual similarity of two images. CLIP loss instead uses the CLIP model to extract features as it learns a powerful shared visual-language embedding. 
%\zj{However, we found these losses have their limitations that ... . Thus,}
Previous works optimize the parameters of the vector sketch with CLIP loss. We found that an extended version of SDS loss can give more diverse sketch synthesis results, and it can be combined with  CLIP loss or LPIPS loss,  providing an additional control term.
%%%Impressed by the results of DreamFusion~\cite{poole_2023_dreamfusion}, we present an extended version of SDS loss for text-to-sketch generation. We find that this loss can encourage the diversity of generated sketches and speed up the synthesis process. 
%and demonstrate that this loss is especially useful in capturing details in generated sketches. 
(2) Explore attention maps embedded in the diffusion model for effective stroke initialization.
% To improve the generation quality and efficiency, we further improve the model from three aspects: (1) Add prior knowledge to stroke initialization. 
The synthesis process can be very time-consuming if the starting point of each stroke is randomly initialized. Therefore, we explore different initialization strategies and present to initialize the curve control points with a fused product of cross-attention maps and self-attention maps in the U-Net~\cite{Dhariwal_ADM_2021} of the diffusion model,
%%%The cross-attention layers regulate the interplay between the spatial layout of the image and each word in the prompt, while the self-attention layer impacts the spatial layout and geometric structure of the generated image~\cite{prompt2prompt_hertz_2023}. 
%We initialize the curve control points by linearly combining the probability distributions of the two attention maps, 
which significantly improves the efficiency compared to the random initialization. 
(3) Introduce opacity property during the optimization of Bézier curves to better mimic the style of human sketches, achieving the effect of heavy and light brushstrokes.
% (3) To achieve the heavy and light painting style as human sketches, we define a sketch as a set of Bézier curves with opacity properties.
%%%Besides improving visual quality, this design implicitly enables the model to adjust the number of strokes. For example, one stroke will disappear once its opacity becomes 1 (100\% transparent). 

% \textcolor{green}{To achieve the heavy and light painting style as human sketches, we define a sketch as a set of Bézier curves with opacity properties and optimize the curves' parameters through a differentiable rasterizer guided by a pretrained diffusion model~\cite{Li_2020_diffvg}.
% Inspired by DreamFusion\cite{poole_2023_dreamfusion}, we propose an input augmentation version of the score distillation sampling (SDS) loss to optimize curve parameters. The SDS loss converts diffusion sampling into an optimization problem that allows the raster image to be represented by a differentiable rasterizer. To accelerate the generation process and improve fidelity, we initialize DiffSketcher with a fused product of cross-attention maps and self-attention maps in the U-Net~\cite{Dhariwal_ADM_2021} used by the diffusion model. The cross-attention layers regulate the interplay between the spatial layout of the image and each word in the prompt, while the self-attention layer impacts the spatial layout and geometric structure of the generated image~\cite{prompt2prompt_hertz_2023}. We initialize the curve control points by linearly combining the probability distributions of the two attention maps, which significantly improves the quality of the final sketch compared to random initialization.}

Our contributions are threefold: 
(1) We propose a text-to-sketch diffusion model dubbed DiffSketcher. To the best of our knowledge, it is the first diffusion model to generate diverse and high-quality vector sketches with different levels of abstraction at the object and scene levels. 
%an algorithm that synthesizes free-hand sketches based on natural language input. 
(2) We present three strategies to improve the generation quality and efficiency, including an extended SDS loss and an attention-enhanced stroke initialization strategy. The opacity property of strokes is considered during synthesis to improve the visual effect further. 
%(2) We introduce an extended version of the SDS loss to optimize curve parameters and leverage a pre-trained latent diffusion model as a prior for optimization. 
(3) We conduct extensive experiments, and the experimental results show that our model outperforms other models in sketch quality and diversity. Thorough insights are provided for future research.
%(3) We improve the generation process and fidelity of DiffSketcher by initializing it with a fused product of cross-attention maps and self-attention maps in the U-Net used by the diffusion model.

\begin{figure}[p]
\centering
\includegraphics[width=\textwidth,height=\textheight,keepaspectratio]{more_results.pdf}
\vspace{-1em}
\caption{
Various free-hand sketches synthesized by DiffSketcher and the corresponding description prompts.
DiffSketcher obtains prior information from LDM~\cite{ldm_2022_Rombach} composite images through score distillation~\cite{poole_2023_dreamfusion} and achieves the same heavy and light drawing styles as human sketches by performing gradient descent on a set of Bézier curves with the opacity property.
Our proposed DiffSketcher allows for varying levels of abstraction while matching its corresponding textual semantics.
In each example, given the same text prompt and two different random seeds, two sketches with a different number of strokes are generated. The red words represent the cross-attention index used to initialize the control points (details about cross-attention are covered in Section~\ref{sec:stroke_init}).
%In the bottom, the same text prompt corresponds to two different random seed and stroke numbers and the red words represent the cross-attention index used to initialize the control points (details about cross-attention are covered in Section~\ref{sec:stroke_init}).
}
\label{fig:more_results}
\end{figure}

\section{Related Work}

\noindent\textbf{Sketch Synthesis.}\quad
Free-hand drawings convey abstract concepts through human visual perception with minimal abstraction. Unlike purely edge-map extraction methods~\cite{canny_1986}, free-hand sketching aims to present sketches that are abstract in terms of structure~\cite{infodrawing_Chan_2022} and semantic interpretation\cite{vinker_2022_clipasso}. Therefore, computational sketching methods that aim to mimic human drawing consider a wide range of sketch representations, ranging from those grounded in the edge map of the input image\cite{xie_2015_holistically,free_li_2017, li2019photo, grayscale_tong_2021,liu2020unsupervised,deepfacedrawing_chen_2020} to those that are more abstract~\cite{sketchRNN_ha_2018a,creative_ge_2021,doodleformer_bhunia_2022, clipdraw_frans_2022, SketchLattice_Qi_2021,vinker_2022_clipasso}, which are normally in vector format. Among the works synthesizing vector sketches, CLIPasso~\cite{vinker_2022_clipasso} and Clipascene~\cite{clipascene_vinker_2022} are conditioned on an input image, while the rest are unconditional. Until now, no prior work has explored synthesizing a sketch based on text. 
%Some work attempts to develop a unified algorithm that can output sketches with various styles~\cite{infodrawing_Chan_2022, neuralstrokes_liu_2021,Portrait_YiLLR_20}. However, few works provide various levels of abstraction~\cite{Portrait_Berger_2013, deep_abstraction_muhammad_2018, vinker_2022_clipasso}.
%In summary, previous research has explored a range of sketch synthesis methods, including those that focus on edge-map extraction and those that aim to produce more abstract sketches. Several works have attempted to develop a unified algorithm that can output sketches with a variety of styles. However, no work has addressed the different levels of abstraction of sketch synthesis based on text prompts. 
%Our work builds upon these prior efforts by proposing an algorithm that synthesizes free-hand sketches based on natural language input, utilizing a pre-trained text-to-image diffusion model and a differentiable rasterizer.

\noindent\textbf{Vector Graphics.}\quad Our work builds upon the differentiable renderer for vector graphics introduced by Li et al~\cite{Li_2020_diffvg}. While image generation methods that operate over vector images traditionally require a vector-based dataset, recent work has shown how differentiable renderers can be used to bypass this limitation~\cite{clipgen_shen_2021,rethinking_kotovenko_2021, evolution_tian_2022, im2vec_reddy_2021, Ma_2022_LIVE}. Furthermore, recent advances in visual text embedding contrastive language-image pre-training (CLIP)\cite{clip_radford_2021} have enabled a number of successful methods for synthesizing sketches, such as CLIPDraw\cite{clipdraw_frans_2022}, StyleCLIPDraw~\cite{Styleclipdraw_schaldenbr_2021}, CLIP-CLOP~\cite{clipclop_mirowski_2022}, and CliPascene~\cite{clipascene_vinker_2022}. A very recent work VectorFusion~\cite{jain2022vectorfusion} combine differentiable renderer with diffusion model for vector graphics generation, \textit{e.g.}, iconography and pixel art. Our proposed algorithm, DiffSketcher, shares a similar idea with VectorFusion, but our focus is generating object- and scene-level sketches from a natural language prompt.

\noindent\textbf{Diffusion Models.} Denoising diffusion probabilistic models (DDPMs)~\cite{diffusion_models_dickstein_2015,EestGrad_song_2019,ddpm_ho_2020,scorebased_song_2021}, particularly those conditioned on text, have shown promising results in text-to-image synthesis. For example, Classifier-Free Guidance (CFG)\cite{classifierfree_2022_ho} has improved sample quality and has been widely used in large-scale diffusion model frameworks, including GLIDE\cite{GLIDE_2022_nichol}, Stable Diffusion~\cite{ldm_2022_Rombach}, DALL·E 2~\cite{DALLE2_2022_ramesh}, and Imagen~\cite{imagen_2022_saharia}. However, the majority of images available in web-scale datasets are rasterized, and this work follows the framework of \textit{synthesis through optimization}, in which images are generated through evaluation-time optimization against a given metric. Our proposed algorithm, \textit{DiffSketcher}, uses a pre-trained text-to-image diffusion model to synthesize free-hand sketches from natural language input. This is achieved by transferring image synthesis prior information into a differentiable renderer.

\section{Preliminaries}

\subsection{Diffusion Models}

In this section, we provide a concise overview of diffusion models, which are a class of generative models that utilize latent variables to gradually transform a sample from a noise distribution to a target data distribution~\cite{diffusion_models_dickstein_2015,ddpm_ho_2020}.
Diffusion models consist of two components: a forward process $q$ and a reverse process or generative model $p$. The forward process, which is typically modeled as a Gaussian distribution, gradually removes structure from the input data $x$ by adding noise over time. The reverse process, on the other hand, adds structure to the noise starting from a latent variable $\mathbf{z}_t$.
Specifically, the generative model is trained to slowly add structure starting from random noise $p(\mathbf{z}_T) = \mathcal{N}(\mathbf{0}, \mathbf{I})$ with transitions $p_{\phi}(\mathbf{z}_{t-1}|\mathbf{z}_{t})$.
%To compute the marginal distribution of the latent variables at timestep $t$ given an initial datapoint $x$, we integrate out intermediate timesteps as follows:
$q(\mathbf{z}_t | \mathbf{x}) = \mathcal{N}(\alpha_t \mathbf{x}, \sigma^2_t \mathbf{I})$.
%Here, $\alpha_t$ and $\sigma_t$ are coefficients chosen to ensure that the marginal distribution $q(\mathbf{z}_t)$ is close to the data distribution at the beginning of the process ($\sigma_T \approx 0$) and close to a Gaussian distribution at the end of the forward process ($\sigma_T \approx 1$), with $\alpha^2_t = 1 - \sigma^2_t$ chosen to preserve the variance~\cite{snr_kingma_2021on, ddpm_ho_2020}.
%Finally, the marginals obtained by integrating out the data density $q(\mathbf{x})$ correspond to smoothed versions of the data distribution $q(\mathbf{z}_t) = \int q(\mathbf{z}_t|\mathbf{x}) q(\mathbf{x}) d\mathbf{x}$.
Training the diffusion model with a (weighted) evidence lower bound (ELBO) simplifies to a weighted denoising score matching objective for parameters $\phi$~\cite{ddpm_ho_2020}
\begin{equation}
\mathcal{L}_{\mathrm{Diff}}(\phi, \mathbf{x})=\mathbb{E}_{t \sim \mathcal{U}(0,1), \epsilon \sim \mathcal{N}(\mathbf{0}, \mathbf{I})}
\left[
w(t)\left\|\epsilon_{\phi}\left(\alpha_{t} \mathbf{x}+\sigma_{t} \epsilon ; t\right)-\epsilon\right\|_{2}^{2}
\right]
\end{equation}
where $w(t)$ is a weighting function that depends on the timestep $t$.
Our work builds on text-to-image latent diffusion model (LDM) that learn $\epsilon_{\phi}(\mathbf{z}_t;t,y)$ conditioned on text embeddings $y$~\cite{ldm_2022_Rombach}.
LDM uses classifier-free guidance(CFG)~\cite{classifierfree_2022_ho}, which jointly learns an unconditional model to enable higher quality generation via a guidance scale parameter $\omega: \hat{\epsilon}_{\phi} (\mathbf{z}_t;y,t) = (1+w) \epsilon_{\phi}(\mathbf{z}_t;y,t) - w\epsilon_{\phi}(\mathbf{z}_t;t)$ ($\hat{\epsilon}_{\phi}$ denotes the guided version of the noise prediction). CFG alters the score function to prefer regions where the ratio of the conditional density to the unconditional density is large. In practice, setting $w > 0$ improves sample fidelity at the cost of diversity.

\subsection{Score Distillation Sampling}

Many 3D generative approaches use a frozen image-text joint embedding model (\textit{e.g.} CLIP) and an optimization-based approach to train a Neural Radiance Fields (NeRF)~\cite{nerf_mildenhall_2021}.
Such models can be specified as a differentiable image parameterization (DIP)~\cite{differentiable_mordvintsev_2018}, where a differentiable generator $g$ transforms parameters $\theta$ to create an image $\mathbf{x} = g(\theta)$.
In DreamFusion~\cite{poole_2023_dreamfusion}, $\theta$ be parameters of a 3D volume and $g$ is a volumetric renderer.
To learn these parameters, DreamFusion proposed \textit{score distillation sampling} (SDS) loss that can be applied to Imagen~\cite{imagen_2022_saharia}:
\begin{equation}
\nabla_{\theta} \mathcal{L}_{\mathrm{SDS}} (\phi, \mathbf{x}=g(\theta)) \triangleq \mathbb{E}_{t,\epsilon} \left[ 
\omega(t) (\hat{\epsilon}_{\phi}(\mathbf{z};y,t) - \epsilon) 
\frac{\partial \mathbf{x}}{\partial \theta}
\right]
\end{equation}
where the constant $\alpha_t \mathbf{I} = \partial \mathbf{z}_t / \partial \mathbf{x}$ absorbed into $w(t)$, and use the classifier-free-guided $\hat{\epsilon}_{\phi}$.
In practice, SDS gives access to loss gradients, not a scalar loss. Their proposed SDS loss provides a way to assess the similarity between an image and a caption:
\begin{equation}
\nabla_{\theta} \mathcal{L}_{\mathrm{SDS}}(\phi, \mathbf{x}=g(\theta))
=
\nabla_{\theta} \mathbb{E}_{t} \left[
\sigma_{t} / \alpha_{t} w(t) 
\mathrm{KL}
\left(
q\left(\mathbf{z}_{t} | g(\theta) ; y, t\right)|p_{\phi}\left(\mathbf{z}_{t} ; y, t\right)
\right)
\right]
\end{equation}
where $p_{\phi}$ is the distribution learned by the frozen, pretrained Imagen model. $q$ is a unimodal Gaussian distribution centered at a learned mean image $g(\theta)$.
DreamFusion~\cite{poole_2023_dreamfusion} proposed an approach to use a pretrained pixel-space text-to-image diffusion model (Imagen~\cite{imagen_2022_saharia}) as a loss function. However, diffusion models trained on pixels have traditionally been used to sample only pixels.
%We are not interested in sampling pixels; we instead want to create what look like good sketches that match the text prompts when rendered from a set of vector strokes. 
We want to create what look like good sketches that match the text prompts when rendered from a set of vector strokes.
Such models can be specified as a differentiable image parameterization, where a differentiable rasterizer $\mathcal{R}$ transforms parameters $\theta$ to create a sketch $\mathcal{S} = \mathcal{R}(\theta)$.

Inspired by DreamFusion~\cite{poole_2023_dreamfusion}, we extend \textit{score distillation sampling} (SDS) loss to use a pretrained latent diffusion model~\cite{ldm_2022_Rombach} as a prior for optimizing curve parameters.
Intuitively, score distillation converts diffusion sampling into an optimization problem that allows the raster image to be represented by a differentiable rasterizer.

\subsection{Differentiable Rasterizer}

Li \textit{et al.}~\cite{Li_2020_diffvg} introduce a differentiable rasterizer $\mathcal{R}$ that bridges the vector graphics and raster image domains. 
A raster image is a 2D grid sampling over the space of the vector graphics scene $f(x,y;\Theta)$, where $\Theta$ contains the curve parameters, \textit{e.g.,} coordinates of Bézier control points, opacity, line thickness. Given a 2D location $(x,y) \in \mathbb{R}^2$, they first find all the filled curves and strokes that overlap with the location. Then sort them with a user-specified order and compute the color using alpha blending~\cite{digital_images_porter_1984}. It enables the gradients of the rasterized curves (image) to be backpropagated to the curve parameters.

\section{Methodology}

In this section, we present our method for generating sketches using pre-trained text-to-image diffusion models.
Let $\mathcal{S}$ be a sketch that was rendered by a differentiable rendering~\cite{Li_2020_diffvg} $\mathcal{R}$ using the text prompt $\mathcal{P}$. Our goal is to optimize a set of parametric strokes to automatically generate a vector sketch that matches the description of the text prompt.

Our pipeline is illustrated in Figure~\ref{fig:pipeline}. Given a text prompt $y$ of the desired subject and a set of control points, we synthesize the corresponding sketch $\mathcal{S}$ while matching the semantic attributes of the text prompt. To initialize the locations of the strokes, we extract and fuse the attention map of the U-Net~\cite{Dhariwal_ADM_2021} used by the latent diffusion model~\cite{ldm_2022_Rombach}. Our key observation is that the structure and appearance of the generated image depend on the interaction between the pixels and the text embedding through the diffusion process~\cite{prompt2prompt_hertz_2023}. We provide more details on strokes initialization in Section~\ref{sec:stroke_init}.
%----------------Figure 3----------------
\begin{wrapfigure}{r}{0.5\textwidth}
% \vskip -0.1in
\centering
\includegraphics[width=0.3\textheight]{pipeline.pdf}
\vspace{-0.1in}
\caption{
\footnotesize Pipeline overview. DiffSketcher accepts a set of control points and text prompts as input to generate a hand-drawn sketch.
}
\label{fig:pipeline}
\vskip -0.15in
\end{wrapfigure}
%--------------------------------
As shown in Figure~\ref{fig:overview}, in each step of the optimization, we feed the stroke parameters to a differentiable rasterizer $\mathcal{R}$ to produce the raster sketch.
We optimize over parameters $\theta$ such that $\mathcal{R} = \mathcal{S}(\theta)$ is close to a sample from the frozen latent diffusion model~\cite{ldm_2022_Rombach}. 
To perform this optimization, we use an input augmentation version of the SDS loss function, where plausible images have low loss and implausible images have a high loss. 
This results in the raster sketch always being coherent with the text prompt. 
The resulting sketch, as well as the sample from the frozen latent diffusion model, then defines a joint semantic and perceptual loss. 
We backpropagate the loss through the differentiable rasterizer $\mathcal{R}$ and update the control points and opacity of strokes directly at each step until convergence of the loss function.

\subsection{Synthesis Through Optimization}

We define a sketch as a set of $n$ strokes $\{s_1, \dots, s_n\}$ with the opacity attribute placed on a white background. 
To represent each stroke, we use a two-dimensional Bézier curve with four control points $s_i = \{p_i^j\}_{j=1}^4 = \{(x_i,y_i)^j\}_{j=1}^4$ and one opacity attribute $o_i$.
We incorporate the opacity of the strokes into the optimization process and use DiffSketcher semantics understanding to achieve a human-like \textit{heavy and light sketch style}.
The parameters of the strokes are fed to a differentiable rasterizer $\mathcal{R}$, which forms the raster sketch $\mathcal{S} = \mathcal{R} ((s_1,o_1) \dots, (s_n, o_n)) = \mathcal{R} ((\{p_1^j\}_{j=1}^4, o_1), \dots, (\{p_1^j\}_{j=1}^4, o_n))$. For simplicity, we define the parameter in $\mathcal{R}$ as $\theta$.

\subsubsection{Vanilla Version: Fidelity to Generated Image}
\label{sec:vanilla_version}
We start with a two-stage pipeline: First, we sample an image from the latent diffusion model~\cite{ldm_2022_Rombach} using a text prompt. Next, we optimize the control points to obtain a sketch that is consistent with the text prompt. 
To preserve the fidelity of the generated sample, we incorporate a \textbf{J}oint \textbf{V}isual \textbf{S}emantic and \textbf{P}erceptual (JVSP) loss to optimize the similarity of the synthesized sketches and the instance sampled from the frozen latent diffusion model~\cite{ldm_2022_Rombach}.
We leverage the VAE~\cite{taming_2021_esser} decoder $\mathcal{D}$ to get the RGB pixel representation of $\hat{\mathcal{S}}_{\phi}(\mathbf{z}_t|y;t)$.
Then, we jointly use the LPIPS~\cite{lpips_zhang_2018} and CLIP visual encoders~\cite{clip_radford_2021,vinker_2022_clipasso} as depth structural similarity and visual semantic similarity metrics, respectively.
Specifically, we use the following loss function:
\begin{equation}
\mathcal{L}_{\mathrm{JVSP}} = 
\mathcal{L}_{\mathrm{LPIPS}}(\mathcal{D}(\hat{\mathcal{S}}_{\phi}(\mathbf{z}_t|y;t)), \mathcal{R}(\theta)) + 
\sum_l \left\| \mathrm{CLIP}_l (\mathcal{D}(\hat{\mathcal{S}}_{\phi}(\mathbf{z}_t|y;t))) - \mathrm{CLIP}_l (\mathcal{R}(\theta)) \right\|
\end{equation}
The $\mathcal{L}_{\mathrm{JVSP}}$ loss function encourages the synthesized sketch to match the underlying semantics and perceptual details of the image sampled from LDM, leading to more realistic and visually appealing results.
While vectorizing a rasterized diffusion sample is lossy, ASDS loss can either finetune the results or optimize random control points from scratch to sample a sketch that is consistent with the text prompt.
In the Sec~\ref{sec:ASDS_loss}, we introduce the ASDS loss to match the text prompt.

\subsubsection{Augmentation SDS Loss: Fidelity to Text Prompt}
\label{sec:ASDS_loss}
To synthesize a vector sketch that matches a given text prompt, we directly optimize the parameters $\theta$ of the differentiable rasterizer $\mathcal{R}$ that produces the raster sketch $\mathcal{S}$.
We propose an input augmentation version of the SDS loss function to perform this optimization, which encourages plausible images to have low loss and implausible images to have a high loss. 
Given a raster sketch $\tilde{\mathcal{S}} \in \mathbb{R}^{H \times W \times 3}$, we combine \textit{RandomPerspective}, \textit{RandomResizedCrop} and \textit{RandomAdjustSharpness} to get a data augmentation version of $\tilde{\mathcal{S}_a} \in \mathbb{R}^{512 \times 512 \times 3}$.
Then, the LDM uses a VAE encoder~\cite{taming_2021_esser} to encode $\tilde{\mathcal{S}_a}$ into a latent representation $\mathbf{z} = \mathcal{E}(\mathbf{\tilde{\mathcal{S}}_a})$, where $\mathbf{z} \in \mathbb{R}^{(H / f) \times (W / f) \times 4}$ and $f$ is the encoder downsample factor.
Then, we use the following ASDS loss function:
\begin{equation}
\nabla_{\theta} \mathcal{L}_{\mathrm{ASDS}} (\phi, \mathcal{S} = \mathcal{R}(\theta)) \triangleq 
\mathbb{E}_{t,\epsilon,a} 
\left[ 
w(t) (\hat{\epsilon}_{\phi}(\mathbf{z}_t;y,t) - \epsilon) 
\frac{\partial \mathbf{z}}{\partial \tilde{\mathcal{S}}_a}
\frac{\partial \tilde{\mathcal{S}}_a}{\partial \theta}
\right]
\end{equation}
Where the weighting function $w(t)$ is a hyper-parameter. And we sample $t \sim \mathcal{U}(0.05, 0.95)$, avoiding very high and low noise levels due to numerical instabilities. Like DreamFusion~\cite{poole_2023_dreamfusion}, we set $\omega = 100$ for classifier-free guidance and higher guidance weights give improved sample quality.
Intuitively, this loss perturbs $\tilde{\mathcal{S}}_a$ with a random amount of noise corresponding to the timestep $t$, and estimates an update direction that follows the score function of the diffusion model to move to a higher density region.
The ASDS loss function encourages the synthesized sketch to match the given text prompt, while also preserving the style and structure of the original sketch.
At each iteration, we backpropagate the loss through the differentiable rasterizer and update the control points and opacity of strokes directly until convergence of the loss function.

\noindent\textbf{Loss objectives of our final model.} To further enhance the quality of the synthesized sketches, we incorporate the JVSP and ASDS losses. Specifically, we first obtain the initial results with the JVSP loss, then we fine-tune the differentiable rasterizer together with the SDS loss. We found such design can achieve the best performance. The ASDS loss predicts the gradient update direction directly, thus its loss value does not require weight balancing. In the JVSP loss, we set the weight of the LPIPS item to 0.2, and the weight of the CLIP visual item to 1. To compute the L2 distance between intermediate level activations of CLIP, we follow the CLIPasso~\cite{clipascene_vinker_2022} method and use layers 3 and 4 of the ResNet101 CLIP model.
As shown in Figure~\ref{fig:ablation}, we compare the performance of the vanilla version with that of the version using only the ASDS loss, and the proposed final version. Our experiments show that the final version improves both the generation quality and efficiency.

\begin{figure*}[t]
\vspace{-1em}
\centering
\includegraphics[width=1\linewidth]{method_main.pdf}
\vspace{-1.5em}
\caption{
Optimization overview. To synthesize a sketch that matches the given text prompt, we optimize the parameters of the differentiable rasterizer $\mathcal{R}$ that produces the raster sketch $\mathcal{S}$, such that the resulting sketch is close to a sample from the frozen latent diffusion model (the blue part of the picture). 
Since the diffusion model directly predicts the update direction, we do not need to backpropagate through the diffusion model; the model simply acts like an efficient, frozen critic that predicts image-space edits.
}
\vspace{-1em}
\label{fig:overview}
\end{figure*}

\subsection{Joint Attention-based Stroke Initialization}
\label{sec:stroke_init}
\begin{figure*}[t]
\vspace{-1em}
\centering
\includegraphics[height=0.25\textheight]{attn_init.pdf}
\vspace{-0.5em}
\caption{
Strokes Initialization. The blue part of the figure represents the UNet in the LDM, which has two types of attention mechanisms: self-attention and cross-attention. 
The yellow and green parts respectively depict the visualization results of the cross-attention and self-attention.
The gray part shows how the initial strokes are generated using a fused attention map. The dashed box represents the attention fusion, which is composed of the mean of the self-attention map and the cross-attention map corresponding to the $5$-th text prompt token (``Tower''). We start at $1$-th token, because $0$-th token is taken up by the CLIP starting token.
}
\label{fig:attn_init}
\vspace{-1em}
\end{figure*}
The highly non-convex nature of the ASDS loss function makes the optimization process susceptible to initialization, especially in multi-instance scenarios where strokes must be carefully placed to emphasize the overall semantics conveyed by free-hand sketching.
We improve the convergence towards semantic depictions by initializing the curve control points based on the attention map of the text prompt conditional latent diffusion model. 
The UNet in the LDM has two types of attention mechanisms (the blue part), self-attention and cross-attention, as shown in Figure~\ref{fig:attn_init} where the yellow part represents the cross-attention visualization result, and the green part represents the self-attention visualization result.
The structure and appearance of the image generated by the LDM~\cite{ldm_2022_Rombach} depend on the interaction between the pixels to the text embedding through the diffusion process~\cite{prompt2prompt_hertz_2023}, which is manifested in the cross-attention layer~\cite{attention_vaswani_2017}.
The visualization results indicate that the cross-attention layers control the relationship between the spatial layout of the image and each word in the prompt, while the self-attention layer affects the spatial layout and geometric structure of the generated image. 
With this observation, we linearly combine the probability distributions of the two attention maps to initialize the curve control points. Specifically, we pick a map in the cross-attention maps based on a text token and combine it with the averaged self-attention map. 
This process is formalized as $\mathrm{FinalAttn} = \lambda * \mathrm{CrossAttn}_i + (1-\lambda) * \mathrm{Mean}(\mathrm{SelfAttn})$, where $\lambda$ is the control coefficient and $i$ indicates the $i$-th token in the text prompt.
Finally, we normalize the fused attention map using softmax and use it as a distribution map to sample $n$ positions for the first control point $p_i^1$ of each Bézier curve. The other three control points ($p_i^2$,$p_i^3$,$p_i^4$) are sampled within a small radius (0.05 of image size) around $p_i^1$ to define the initial set of Bézier curves $\{ \{p_i^j\}_{j=1}^4 \}_{i=1}^n$. Empirical results show that our attentional fusion-based initialization contributes significantly to the quality and rendering speed of the final sketch in comparison to the random initialization.


% The overall synthesis through optimization process is summarized in Algorithm~\ref{algo:ours_method}.
% \begin{algorithm}[h]
% \caption{DiffSketch}
% \label{algo:ours_method}
% \begin{algorithmic}[1]
% \Require the text prompt $y$; Iteration Count $\mathbf{I}$; Curve Count $N$; Augment Size $\mathbf{A}$; 
% \State Load pre-trained text-to-image diffusion model; \Comment{the start point}
% \State Initialize DiffSketch with cross-attention maps and self-attention maps in U-Net;
% \For{$i=I$ {\bfseries to} $1$}  \Comment{optimizat1}
% \State $s \xrightarrow{} i h$.
% \EndFor
% \State $\boldsymbol{y}_0 \xleftarrow{} \boldsymbol{y}$
% \State \textbf{return} $\mathcal{S}$
% \end{algorithmic}
% \end{algorithm}


\section{Results}
\begin{figure}[h]
\vspace{-1.5em}
\centering 
\includegraphics[width=1\linewidth]{compare.pdf}
\vspace{-1.5em}
\caption{
Comparison with Existing Methods.
}
\vspace{-1.5em}
\label{fig:compare}
\end{figure}
\subsection{Qualitative Evaluation}
As shown in Figure ~\ref{fig:more_results}, we demonstrate that our approach offers the ability to produce object-level and scene-level sketches based on a textual prompt, with the flexibility to manipulate the level of abstraction through stroke count. It is effective in generating accurate sketches regardless of prompt complexity, including simple or non-existent objects, iconic constructions, and detailed scenes.  And the utilization of the robust prior of stable diffusion allows for a favorable initialization, promoting the production of high-quality sketches with significantly fewer iterations.

\subsection{Comparison with Existing Methods}
Figure~\ref{fig:compare} qualitatively compared our DiffSketcher with those of other CLIP-based~\cite{vinker_2022_clipasso} and edge extraction-based methods~\cite{canny_1986}.
The Canny edge extraction algorithm extracts an excessive number of edges and produces untidy sketches, as observed in the first example in Figure~\ref{fig:compare}. CLIPasso uses visual distance metrics to guide gradient-based optimization. On the task of drawing scene-level sketches, CLIPasso can only draw part of the foreground, and the background part is missing. DiffSketcher produces much cleaner vector sketches than baselines because we incorporate a generative prior of image appearance.

\begin{figure}[t]
\centering 
\vspace{-0.5em}
\includegraphics[width=1\linewidth]{ablation.pdf}
\vspace{-1em}
\caption{
Qualitative results of ablation study. \textbf{Top}: two initialization strategies used by DiffSketcher (\textit{i.e.}, LDM Attention) and CLIPasso~\cite{vinker_2022_clipasso} (\textit{i.e.}, CAM~\cite{gradcam_selvaraju_2017} Saliency map). A result of sampling from LDM . The \textbf{2nd} and \textbf{3rd} rows: effect comparison of the two initialization strategies on the rate of convergence. The \textbf{4th} and \textbf{5th} rows: The effect of JVSP (Section~\ref{sec:vanilla_version}) and ASDS loss (Section~\ref{sec:ASDS_loss}). The \textbf{6th} and \textbf{7th} rows: ASDS loss can lead to more diversity when used with random initialization. Text prompt used in the example: ``Astronaut on Asteroid, galaxy background''.
}
\label{fig:ablation}
\end{figure}

\subsection{Quantitative Evaluation}
Evaluating text-to-sketch synthesis is challenging due to the absence of ground truth sketches. Therefore, we focus on three indicators: consistency between the generated sketch and text prompt, the aesthetic quality of the sketch, and the recognition accuracy of the sketch.
To measure the consistency between the generated sketch and input text, we calculate the mean of cosine similarity of CLIP embeddings for the generated sketches and the text captions used to generate them. Our method achieves a cosine similarity of $0.3494$, which is higher than the $0.328$ achieved by the Canny algorithm and the $0.3075$ achieved by CLIPasso.
Aesthetic appeal is subjective, and beauty is a personal experience and preference. However, the beauty of a hand-drawn sketch can be described in terms of various factors, such as line quality, texture, material, composition, style, and expressiveness. As an evaluation of our proposed method, we use CLIP-based aesthetic indicators~\cite{aesthetic_predictor} to calculate the aesthetic value for samples of multiple categories. Figure~\ref{fig:compare} shows the aesthetic value of various methods in some examples. Our method achieves a mean value of $4.8206$ for the aesthetic value on a large number of examples, which is higher than the $4.3682$ achieved by the Canny algorithm and the $4.0821$ achieved by CLIPasso~\cite{vinker_2022_clipasso}.

\subsection{Ablation Study}
We have conducted a series of ablation studies to show the effect of the proposed initialization strategy and ASDS loss.
As shown in Figure~\ref{fig:ablation}, the initialization mode is critical for non-convex objective function optimization, and the proposed initialization strategy improves sampling quality and efficiency. 
% The LDM's cross attention feature can efficiently activate the relevant regions based on the token areas, while the self attention mechanism can effectively differentiate the foreground from the background down to the pixel level. By combining both cross attention and self attention mechanisms, the initialization area becomes more precise. This allows the renderer to place more initial points within the active region corresponding to the foreground object, resulting in a superior quality sketch with fewer optimization steps required.We used the LDM attention map and saliency map respectively to select the initial points for the renderer, and observed the convergence rate of the generated sketch on our method and clipasso respectively. It is evident from Figure ~\ref{fig:compare} that the linear attention map generated by LDM shows a more precise representation of the foreground object region in comparison to the saliency map. A precise identification of the foreground object region through this method enables the attainment of highly suitable initial points, leading to faster convergence of the system.
% When random initial points are used, the ASDS loss can guide the rendering process to produce morphic sketches. However, without any semantic information, the resulting sketches may lack precision. On the other hand, when initial points are informed by semantic knowledge, the ASDS optimization gains direction and stability, ultimately producing more refined sketches. Moreover, the integration of ASDS losses accelerates the model convergence process while also enhancing the level of detail in the final sketch.



\section{Conclusion}
In this work, we have proposed a novel and effective approach to bridge the gap between natural language and free-hand sketches. By leveraging the power of pretrained text-to-image diffusion models, we have developed a method that can generate diverse and high-quality vector sketches based on textual input, without the need for large-scale datasets of sketches or sketch-text pairs.
%Our experiments demonstrate the effectiveness of DiffSketch, both in terms of the quality and diversity of the generated sketches. 
We have also explored different aspects of the model design, such as the stroke initialization strategies, choice of loss functions, and properties of strokes in the differentiable rasterizer, which provide valuable insights for future research.
%In conclusion, our work presents a significant step towards the ultimate goal of enabling machines to understand and communicate abstract concepts in the same way that humans do. 
With further research and development, our proposed model, DiffSketcher, has the potential to be a valuable tool in various domains, such as design and education.

\medskip

\newpage
\clearpage

\bibliography{bibref}
\bibliographystyle{plain}

%%%%%%%%%%%%%%%%%%%%%%%%%%%%%%%%%%%%%%%%%%%%%%%%%%%%%%%%%%%%


\end{document}